\subsection{Entwickler-Onboarding}
\label{subsec:entwickler-onboarding}

Die README trennt zwei Einstiege. Für Produktentwicklung wird das
Master-Repository geklont, mit \texttt{doctor} geprüft und mit \texttt{init}
ein Profilprojekt erzeugt. Erst im Zielverzeichnis folgen erneut
\texttt{doctor}, dann \texttt{install} und \texttt{run}; Produktcode,
Konfiguration, Tests und Commits gehören ab diesem Punkt in das abgeleitete
Repository. Template-Wartung arbeitet dagegen direkt an Baseline, Profilen
oder Tooling des Masters und führt dort mindestens Diagnose, Installation und
den vollständigen Testweg aus. Diese Trennung verhindert, dass Produktarbeit
versehentlich die gemeinsame Vorlage verändert
\parencite{kleiveist2026template,kleiveist2026profiles}.

Als Grundausstattung dokumentiert das Projekt Git, Python~3.11 oder neuer und
Node.js~20 oder neuer mit npm. Desktopprofile ergänzen Rust Stable. Hinzu
kommen die plattformspezifischen Tauri-Voraussetzungen. Docker/Compose wird erst für
Containerpfade benötigt, ein erreichbares PostgreSQL nur für die optionale
PostgreSQL-Capability und PyGitIndex nur zur Regeneration der
Dokumentationsnavigation. Diese Abstufung entspricht den getrennten
Diagnosewegen und verhindert, optionale Werkzeuge als Blocker jedes Profils zu
behandeln.

Die Navigation ist ebenfalls aufgabenbezogen: Die README beschreibt den
Ersteinstieg, der Tooling Guide Befehle und Diagnose, Project Profiles das
Strukturmodell und Runtime Configuration die Wertepriorität. CI-, Container-
und Release-Dokumente führen zu automatisierten Prüfungen und
Auslieferungsgrenzen weiter. Damit reduziert das Repository die Zahl
unterschiedlicher öffentlicher Setup-Pfade; eine messbare Verkürzung des
Onboardings oder Produktivitätssteigerung folgt daraus noch nicht und bleibt
der Bewertung in Kapitel~\ref{sec:bewertung} vorbehalten
\parencite{kleiveist2026tooling,kleiveist2026configuration,kleiveist2026ci,kleiveist2026release}.
